\documentclass[10pt,a4paper]{amsart}
\usepackage[margin=19mm]{geometry}
\usepackage{amsmath,amssymb,mathtools}
\usepackage[scr=rsfs,cal=euler]{mathalfa}
\usepackage{xfrac}
\usepackage[unicode,hidelinks,bookmarks=false]{hyperref}
\newcommand{\ie}{i.e.\ }
\newcommand{\cf}{cf.\ }
\newcommand{\resp}{respectively\ }
\newcommand{\Subsection}{\subsection}
\providecommand{\nobreakdash}{\nobreak\hbox{-}}
\setlength{\emergencystretch}{2em}
\multlinegap0pt
\usepackage{bm}
\usepackage{bbm}
\usepackage{dsfont}
\usepackage{xspace}
\usepackage{cancel}
\usepackage{mathtools}
\usepackage{amsthm}
\makeatletter
\@ifundefined{theorem}{\newtheorem{theorem}{Theorem}}{}
\@ifundefined{lemma}{\newtheorem{lemma}[theorem]{Lemma}}{}
\makeatother
\makeatletter
\expandafter\ifx\csname proof\endcsname\relax
\newenvironment{proof}[1][Proof]{\textbf{#1.} }{\hfill\rule{0.5em}{0.5em}}
\fi
\makeatother
\title[On Quadrature Surface Free Boundary Problems for the Riemann]{On Quadrature Surface Free Boundary Problems for the Riemannian $p$-Laplacian}
\author{Ababacar Sadikhe Djite, Diaraf Seck}
\date{Source: arXiv:2609.17978v1 (CC BY 4.0)}
\begin{document}
\maketitle

\noindent\emph{Source and licence.} On Quadrature Surface Free Boundary Problems for the Riemannian $p$-Laplacian. Version arXiv:2609.17978v1 (CC BY 4.0), distributed under the Creative Commons Attribution 4.0 International licence, as recorded in the arXiv licence metadata for this exact version and shown on the abstract page https://arxiv.org/abs/2609.17978v1. Full rights notice: licenses/arxiv-2609.17978-CC-BY-4.0.txt.\par\smallskip

\noindent\textit{Editorial scope.} Counted unit: the Theorem whose source label is \texttt{thm:stability-p-laplace} in the article (source0.tex lines 1010--1031, source label \texttt{thm:stability-p-laplace}), delivered together with the article's own complete original proof of it (source0.tex lines 1033--1292, one \texttt{proof} environment of 260 lines). No mathematics is written, repaired or re-derived: statement and proof are the article's own text line for line. Required context retained uncounted: lem:composition-C1-W1p (lines 965--983). Every definition, display and label of those lines is retained unchanged. Omitted from this package and not redistributed: 1--964, 984--1009, 1032--1032, 1293--3134. Nothing inside a retained excerpt is omitted or altered: every retained body is delivered exactly as the article's lines are written, so no omission receipt of any kind accompanies this package. Citations: the retained excerpts carry no citation command at all, so no bibliography entry of the article is transcribed and no reference is redistributed. Reference adaptation: none was needed and none was made; every reference inside the delivered unit resolves inside it, so no unresolved reference or citation is produced. Printed slips kept literal: no printed word, symbol, hypothesis or number of the counted statement or of its proof was altered. Layout and class: the article's own class is \texttt{article} with packages including amssymb, graphicx, amsmath, float, array, comment, fontenc, inputenc; the wrapper is the lane's pinned \texttt{amsart} editorial wrapper at 10pt on A4, which restates the theorem-like environments the retained text needs and adds the article's own macro definitions that the retained text needs (none), with extra wrapper packages (bm, bbm, dsfont, xspace, cancel, mathtools). The delivered numbering is the unit's own. Assets: no retained span contains a figure environment, a graphics-inclusion command, a PDF-page inclusion, a table or a TikZ picture; no image file is copied into this package, the package declares no asset dependency and no image file is redistributed with it. Licence: version arXiv:2609.17978v1 of this article is distributed under the Creative Commons Attribution 4.0 International licence, as recorded in the arXiv licence metadata for this exact version and shown on the abstract page https://arxiv.org/abs/2609.17978v1. A full rights notice is delivered at licenses/arxiv-2609.17978-CC-BY-4.0.txt. Characters: every retained excerpt is pure ASCII, so the delivered engine, XeLaTeX with its default Latin Modern text font, reports no missing character.
\begin{lemma}[Continuity of composition under $C^1$ convergence]
\label{lem:composition-C1-W1p}
Let $(M,g)$ be a compact smooth Riemannian manifold and let $1<p<\infty$. Suppose that $\Phi_j:M\to M$ are $C^1$ diffeomorphisms such that
\[
\Phi_j\to\operatorname{Id}\quad\text{in }C^1(M).
\]
Then, for every $w\in W^{1,p}(M)$,
\[
w\circ\Phi_j\to w\quad\text{strongly in }W^{1,p}(M),
\]
and the same conclusion holds for $\Phi_j^{-1}$. Moreover, the composition
operators induced by $\Phi_j$ and $\Phi_j^{-1}$ are uniformly bounded on
$W^{1,p}(M)$. More generally, for every $1\le q<\infty$ and every
$h\in L^q(M)$,
\[
h\circ\Phi_j\to h\quad	ext{strongly in }L^q(M),
\]
and the same holds for $\Phi_j^{-1}$.
\end{lemma}

\begin{theorem}[Stability of the Riemannian $p$-Laplace problem]
\label{thm:stability-p-laplace}
Let $\Omega_j,\Omega\in\mathcal A_C(M)$ and assume that
\[
\Omega_j\longrightarrow\Omega
\]
in the strong $RC$-$GNP$ topology. Let
$f\in L^{p'}(M)$, where $p'=p/(p-1)$, and let
$u_j=u_{\Omega_j}$ and $u=u_\Omega$ be the corresponding weak
solutions of the Riemannian $p$-Laplace Dirichlet problem, extended by
zero to $M$. Then
\[
u_j\longrightarrow u
\qquad\text{strongly in }W^{1,p}(M).
\]
In particular,
\[
\int_M|\nabla_g u_j|_g^p\,dV_g
\longrightarrow
\int_M|\nabla_g u|_g^p\,dV_g.
\]
\end{theorem}

\begin{proof}
Let $\Phi_j:M\to M$ be the diffeomorphisms associated with the strong
$RC$-$GNP$ convergence, so that
\[
\Phi_j(\Omega)=\Omega_j,
\qquad
\Phi_j\to\operatorname{Id}
\quad\text{in }C^1(M).
\]
Set
\[
g_j:=\Phi_j^*g.
\]
Then
\[
g_j\to g\qquad\text{uniformly on }M,
\]
and, since $\Phi_j\to\operatorname{Id}$ in $C^1$ and the $\Phi_j$ are
diffeomorphisms, the metrics $g_j$ are uniformly equivalent to $g$.
In particular, there exists $c\geq1$, independent of $j$, such that
\[
c^{-1}|\xi|_g\leq |\xi|_{g_j}\leq c|\xi|_g
\]
for every $\xi\in TM$.

Define
\[
v_j=u_j\circ\Phi_j\qquad\text{in }\Omega.
\]
Then
\[
v_j\in W_0^{1,p}(\Omega).
\]
After the change of variables $y=\Phi_j(x)$, the weak formulation on
$\Omega_j$ becomes
\begin{equation}
\int_\Omega
|\nabla_{g_j}v_j|_{g_j}^{p-2}
g_j(\nabla_{g_j}v_j,\nabla_{g_j}\varphi)
\,dV_{g_j}
=
\int_\Omega (f\circ\Phi_j)\varphi\,dV_{g_j}
\label{eq:transported-p-laplace}
\end{equation}
for every $\varphi\in W_0^{1,p}(\Omega)$.

For later use, define the nonlinear form
\[
B_j(w,\varphi)
=
\int_\Omega
|\nabla_{g_j}w|_{g_j}^{p-2}
g_j(\nabla_{g_j}w,\nabla_{g_j}\varphi)
\,dV_{g_j}
\]
and
\[
F_j(\varphi)
=
\int_\Omega(f\circ\Phi_j)\varphi\,dV_{g_j}.
\]
The uniform equivalence of $g_j$ and $g$, together with the uniform
Poincar\'e inequality on the fixed domain $\Omega$, gives
\[
\|v_j\|_{W^{1,p}(\Omega)}\leq C.
\]
Indeed, taking $\varphi=v_j$ in \eqref{eq:transported-p-laplace},
using H\"older's inequality and Poincar\'e's inequality, yields the
bound with $C$ independent of $j$.

Since $W_0^{1,p}(\Omega)$ is reflexive, after extraction of a
subsequence there exists $v\in W_0^{1,p}(\Omega)$ such that
\[
v_j\rightharpoonup v
\qquad\text{weakly in }W_0^{1,p}(\Omega).
\]
Moreover, by the Rellich--Kondrachov theorem,
\[
v_j\longrightarrow v
\qquad\text{strongly in }L^p(\Omega).
\]

We next identify the weak limit. The passage to the limit in the nonlinear term cannot be justified by weak convergence alone. We therefore use the monotonicity of the transported operators.

For an arbitrary $w{ \in }W_0^{1,p}(\Omega)$, monotonicity gives
\[
0\leq B_j(v_j,v_j-w)-B_j(w,v_j-w).
\]
The first term equals $F_j(v_j-w)$ by the transported weak formulation.
By the $L^{p'}$ version of Lemma~\ref{lem:composition-C1-W1p},
\[
f\circ\Phi_j\longrightarrow f
\qquad\text{strongly in }L^{p'}(M).
\]
Together with $v_j\to v$ strongly in $L^p(\Omega)$ and the uniform
convergence of the volume densities $dV_{g_j}$ to $dV_g$, this gives
\[
F_j(v_j-w)\longrightarrow \int_\Omega f(v-w)\,dV_g.
\]
For the second term, $w$ is fixed and the coefficients of the transported
operator converge uniformly. Hence the weak convergence of $v_j$ gives
\[
B_j(w,v_j-w)\longrightarrow
\int_\Omega
|\nabla_g w|_g^{p-2}
 g(\nabla_g w,\nabla_g(v-w))\,dV_g.
\]
It follows that, for every $w\in W_0^{1,p}(\Omega)$,
\[
\int_\Omega
|\nabla_g w|_g^{p-2}
 g(\nabla_g w,\nabla_g(w-v))\,dV_g
\geq
\int_\Omega f(w-v)\,dV_g.
\]
This is the Minty variational inequality for the limiting operator. Taking
$w=v-t\varphi$, with $t>0$ and arbitrary $\varphi\in W_0^{1,p}(\Omega)$,
and then dividing by $t$, we obtain, as $t\downarrow0$,
\[
-\int_\Omega
|\nabla_g(v-t\varphi)|_g^{p-2}
 g(\nabla_g(v-t\varphi),\nabla_g\varphi)\,dV_g
\geq
-\int_\Omega f\varphi\,dV_g.
\]
Using the continuity of the map
$\xi\mapsto|\xi|_g^{p-2}\xi$ from $L^p$ into $L^{p'}$, we obtain
\[
\int_\Omega
|\nabla_g v|_g^{p-2}
 g(\nabla_g v,\nabla_g\varphi)\,dV_g
\leq
\int_\Omega f\varphi\,dV_g.
\]
Repeating the argument with $w=v+t\varphi$ yields the reverse inequality.
Therefore
\[
\int_\Omega
|\nabla_g v|_g^{p-2}
 g(\nabla_g v,\nabla_g\varphi)\,dV_g
=
\int_\Omega f\varphi\,dV_g
\]
for every $\varphi\in W_0^{1,p}(\Omega)$. Thus $v$ is a weak solution of
the limiting Dirichlet problem. By uniqueness,
\[
v=u.
\]
Consequently,
\[
v_j\rightharpoonup u
\qquad\text{weakly in }W_0^{1,p}(\Omega).
\]

It remains to prove strong convergence. Taking $\varphi=v_j-u$ in
\eqref{eq:transported-p-laplace} gives
\[
B_j(v_j,v_j-u)=F_j(v_j-u)\longrightarrow0.
\]
Moreover, since $u$ is fixed and $g_j\to g$ uniformly,
\[
B_j(u,v_j-u)\longrightarrow0
\]
by the weak convergence of $v_j$ to $u$. Hence
\begin{equation}
B_j(v_j,v_j-u)-B_j(u,v_j-u)\longrightarrow0.
\label{eq:monotonicity-limit}
\end{equation}

For $p\ge2$, the uniform monotonicity inequality gives, for some
$c_p>0$ independent of $j$,
\[
\left(
|\xi|_{g_j}^{p-2}\xi-|\eta|_{g_j}^{p-2}\eta
\right)\cdot_{g_j}(\xi-\eta)
\ge c_p|\xi-\eta|_{g_j}^p.
\]
Using \eqref{eq:monotonicity-limit}, we obtain
\[
\|\nabla_{g_j}v_j-\nabla_{g_j}u\|_{L^p(\Omega,g_j)}\longrightarrow0.
\]

For $1<p<2$, we use the standard weighted monotonicity estimate.
The quotient below is understood to be $0$ when $\xi=\eta=0$.
For $1<p<2$, we use the uniform estimate
\[
\left(
|\xi|_{g_j}^{p-2}\xi-|\eta|_{g_j}^{p-2}\eta
\right)\cdot_{g_j}(\xi-\eta)
\ge c_p
\frac{|\xi-\eta|_{g_j}^2}
{(|\xi|_{g_j}+|\eta|_{g_j})^{2-p}}.
\]
It follows from \eqref{eq:monotonicity-limit} that
\[
\int_\Omega
\frac{|\nabla_{g_j}v_j-\nabla_{g_j}u|_{g_j}^2}
{(|\nabla_{g_j}v_j|_{g_j}+|\nabla_{g_j}u|_{g_j})^{2-p}}
\,dV_{g_j}\longrightarrow0.
\]
By H\"older's inequality,
\begin{align*}
\int_\Omega
|\nabla_{g_j}v_j-\nabla_{g_j}u|_{g_j}^p\,dV_{g_j}
&\leq
\left(
\int_\Omega
\frac{|\nabla_{g_j}v_j-\nabla_{g_j}u|_{g_j}^2}
{(|\nabla_{g_j}v_j|_{g_j}+|\nabla_{g_j}u|_{g_j})^{2-p}}
\,dV_{g_j}
\right)^{p/2}
\\
&\quad\times
\left(
\int_\Omega
(|\nabla_{g_j}v_j|_{g_j}+|\nabla_{g_j}u|_{g_j})^p
\,dV_{g_j}
\right)^{(2-p)/2}.
\end{align*}
The second factor is uniformly bounded, so the first one implies
\[
\|\nabla_{g_j}v_j-\nabla_{g_j}u\|_{L^p(\Omega,g_j)}\longrightarrow0.
\]
Thus, for every $1<p<\infty$,
\[
v_j\longrightarrow u
\qquad\text{strongly in }W_0^{1,p}(\Omega).
\]

Let $\widetilde v_j$ and $\widetilde u$ denote the zero extensions of
$v_j$ and $u$ from $\Omega$ to $M$. Since
$v_j\to u$ strongly in $W_0^{1,p}(\Omega)$, we have
\[
\widetilde v_j\to\widetilde u\quad\text{strongly in }W^{1,p}(M).
\]
By construction, $\widetilde v_j=\widetilde u_j\circ\Phi_j$ on $M$,
where $\widetilde u_j$ is the zero extension of $u_j$. Hence
$\widetilde u_j=\widetilde v_j\circ\Phi_j^{-1}$. Therefore, by
Lemma~\ref{lem:composition-C1-W1p},
\begin{align*}
\|\widetilde u_j-\widetilde u\|_{W^{1,p}(M)}
&\le C\|\widetilde v_j-\widetilde u\|_{W^{1,p}(M)}
+\|\widetilde u\circ\Phi_j^{-1}-\widetilde u\|_{W^{1,p}(M)}
\longrightarrow0.
\end{align*}
Thus $u_j\to u$ strongly in $W^{1,p}(M)$.
In particular,
\[
\|\nabla_g u_j\|_{L^p(M)}
\longrightarrow
\|\nabla_g u\|_{L^p(M)},
\]
and hence
\[
\int_M|\nabla_g u_j|_g^p\,dV_g
\longrightarrow
\int_M|\nabla_g u|_g^p\,dV_g.
\]

\end{proof}

\end{document}
